\section{未锚定账本事件的叙事补充：jin}

\label{sec:anchor-batch-two-jin}

本组为尚无正文跳转目标的账本事件补入锚点。每一记录置于战争、行政、环境、迁徙与地方社会的关系中，并按来源限度约束叙事。

\subsection{制度、区域与材料边界}

\eventanchor{JIN-1179-REGNAL-YEAR}{1179年·金以金世宗在位第18年作为诏令、赋役与外交文书的纪年框架}账本条目记录金在1179年的“金以金世宗在位第18年作为诏令、赋役与外交文书的纪年框架”。其背景是新岁颁历、年号和纪年是中枢与地方行政沟通的共同前提，后续影响为为同年政令、战事和地方材料提供日期锚点，不能单独推知具体政策执行。政权、军队或制度的变化须经由道路、文书、资源和地方中介才会落实，城市、乡村、边地和流动人群不可能在同一时间承受相同结果。

本条使用《金史》〈本纪〉、〈交聘表〉及相关志；《宋史》与《建炎以来系年要录》相关条，并标示“纪年框架可靠；该记录仅确证政权纪年连续，年内具体事项须由本年条另证”。这些材料能够钉住年序、命令、战役或局部地点，却不能直接推出人口、税收、身份和社会动机的总量。应区分诏令与实施、条约与边境生活、军报与伤亡统计，也不将官府的族群或行政称谓写成固定的社会本质。

从区域史角度看，该事件是资源、信息与权力重新连接的节点。不同地方会因市场、交通、宗教组织、家庭策略和环境条件而产生不同反应；材料的沉默限制我们对未被记录者所能作出的判断。

\eventanchor{JIN-1180-REGNAL-YEAR}{1180年·金以金世宗在位第19年作为诏令、赋役与外交文书的纪年框架}账本条目记录金在1180年的“金以金世宗在位第19年作为诏令、赋役与外交文书的纪年框架”。其背景是新岁颁历、年号和纪年是中枢与地方行政沟通的共同前提，后续影响为为同年政令、战事和地方材料提供日期锚点，不能单独推知具体政策执行。政权、军队或制度的变化须经由道路、文书、资源和地方中介才会落实，城市、乡村、边地和流动人群不可能在同一时间承受相同结果。

本条使用《金史》〈本纪〉、〈交聘表〉及相关志；《宋史》与《建炎以来系年要录》相关条，并标示“纪年框架可靠；该记录仅确证政权纪年连续，年内具体事项须由本年条另证”。这些材料能够钉住年序、命令、战役或局部地点，却不能直接推出人口、税收、身份和社会动机的总量。应区分诏令与实施、条约与边境生活、军报与伤亡统计，也不将官府的族群或行政称谓写成固定的社会本质。

从区域史角度看，该事件是资源、信息与权力重新连接的节点。不同地方会因市场、交通、宗教组织、家庭策略和环境条件而产生不同反应；材料的沉默限制我们对未被记录者所能作出的判断。

\eventanchor{JIN-1181-REGNAL-YEAR}{1181年·金以金世宗在位第20年作为诏令、赋役与外交文书的纪年框架}账本条目记录金在1181年的“金以金世宗在位第20年作为诏令、赋役与外交文书的纪年框架”。其背景是新岁颁历、年号和纪年是中枢与地方行政沟通的共同前提，后续影响为为同年政令、战事和地方材料提供日期锚点，不能单独推知具体政策执行。政权、军队或制度的变化须经由道路、文书、资源和地方中介才会落实，城市、乡村、边地和流动人群不可能在同一时间承受相同结果。

本条使用《金史》〈本纪〉、〈交聘表〉及相关志；《宋史》与《建炎以来系年要录》相关条，并标示“纪年框架可靠；该记录仅确证政权纪年连续，年内具体事项须由本年条另证”。这些材料能够钉住年序、命令、战役或局部地点，却不能直接推出人口、税收、身份和社会动机的总量。应区分诏令与实施、条约与边境生活、军报与伤亡统计，也不将官府的族群或行政称谓写成固定的社会本质。

从区域史角度看，该事件是资源、信息与权力重新连接的节点。不同地方会因市场、交通、宗教组织、家庭策略和环境条件而产生不同反应；材料的沉默限制我们对未被记录者所能作出的判断。

\eventanchor{JIN-1182-REGNAL-YEAR}{1182年·金以金世宗在位第21年作为诏令、赋役与外交文书的纪年框架}账本条目记录金在1182年的“金以金世宗在位第21年作为诏令、赋役与外交文书的纪年框架”。其背景是新岁颁历、年号和纪年是中枢与地方行政沟通的共同前提，后续影响为为同年政令、战事和地方材料提供日期锚点，不能单独推知具体政策执行。政权、军队或制度的变化须经由道路、文书、资源和地方中介才会落实，城市、乡村、边地和流动人群不可能在同一时间承受相同结果。

本条使用《金史》〈本纪〉、〈交聘表〉及相关志；《宋史》与《建炎以来系年要录》相关条，并标示“纪年框架可靠；该记录仅确证政权纪年连续，年内具体事项须由本年条另证”。这些材料能够钉住年序、命令、战役或局部地点，却不能直接推出人口、税收、身份和社会动机的总量。应区分诏令与实施、条约与边境生活、军报与伤亡统计，也不将官府的族群或行政称谓写成固定的社会本质。

从区域史角度看，该事件是资源、信息与权力重新连接的节点。不同地方会因市场、交通、宗教组织、家庭策略和环境条件而产生不同反应；材料的沉默限制我们对未被记录者所能作出的判断。

\eventanchor{JIN-1183-REGNAL-YEAR}{1183年·金以金世宗在位第22年作为诏令、赋役与外交文书的纪年框架}账本条目记录金在1183年的“金以金世宗在位第22年作为诏令、赋役与外交文书的纪年框架”。其背景是新岁颁历、年号和纪年是中枢与地方行政沟通的共同前提，后续影响为为同年政令、战事和地方材料提供日期锚点，不能单独推知具体政策执行。政权、军队或制度的变化须经由道路、文书、资源和地方中介才会落实，城市、乡村、边地和流动人群不可能在同一时间承受相同结果。

本条使用《金史》〈本纪〉、〈交聘表〉及相关志；《宋史》与《建炎以来系年要录》相关条，并标示“纪年框架可靠；该记录仅确证政权纪年连续，年内具体事项须由本年条另证”。这些材料能够钉住年序、命令、战役或局部地点，却不能直接推出人口、税收、身份和社会动机的总量。应区分诏令与实施、条约与边境生活、军报与伤亡统计，也不将官府的族群或行政称谓写成固定的社会本质。

从区域史角度看，该事件是资源、信息与权力重新连接的节点。不同地方会因市场、交通、宗教组织、家庭策略和环境条件而产生不同反应；材料的沉默限制我们对未被记录者所能作出的判断。

\eventanchor{JIN-1184-REGNAL-YEAR}{1184年·金以金世宗在位第23年作为诏令、赋役与外交文书的纪年框架}账本条目记录金在1184年的“金以金世宗在位第23年作为诏令、赋役与外交文书的纪年框架”。其背景是新岁颁历、年号和纪年是中枢与地方行政沟通的共同前提，后续影响为为同年政令、战事和地方材料提供日期锚点，不能单独推知具体政策执行。政权、军队或制度的变化须经由道路、文书、资源和地方中介才会落实，城市、乡村、边地和流动人群不可能在同一时间承受相同结果。

本条使用《金史》〈本纪〉、〈交聘表〉及相关志；《宋史》与《建炎以来系年要录》相关条，并标示“纪年框架可靠；该记录仅确证政权纪年连续，年内具体事项须由本年条另证”。这些材料能够钉住年序、命令、战役或局部地点，却不能直接推出人口、税收、身份和社会动机的总量。应区分诏令与实施、条约与边境生活、军报与伤亡统计，也不将官府的族群或行政称谓写成固定的社会本质。

从区域史角度看，该事件是资源、信息与权力重新连接的节点。不同地方会因市场、交通、宗教组织、家庭策略和环境条件而产生不同反应；材料的沉默限制我们对未被记录者所能作出的判断。

\eventanchor{JIN-1185-REGNAL-YEAR}{1185年·金以金世宗在位第24年作为诏令、赋役与外交文书的纪年框架}账本条目记录金在1185年的“金以金世宗在位第24年作为诏令、赋役与外交文书的纪年框架”。其背景是新岁颁历、年号和纪年是中枢与地方行政沟通的共同前提，后续影响为为同年政令、战事和地方材料提供日期锚点，不能单独推知具体政策执行。政权、军队或制度的变化须经由道路、文书、资源和地方中介才会落实，城市、乡村、边地和流动人群不可能在同一时间承受相同结果。

本条使用《金史》〈本纪〉、〈交聘表〉及相关志；《宋史》与《建炎以来系年要录》相关条，并标示“纪年框架可靠；该记录仅确证政权纪年连续，年内具体事项须由本年条另证”。这些材料能够钉住年序、命令、战役或局部地点，却不能直接推出人口、税收、身份和社会动机的总量。应区分诏令与实施、条约与边境生活、军报与伤亡统计，也不将官府的族群或行政称谓写成固定的社会本质。

从区域史角度看，该事件是资源、信息与权力重新连接的节点。不同地方会因市场、交通、宗教组织、家庭策略和环境条件而产生不同反应；材料的沉默限制我们对未被记录者所能作出的判断。

\eventanchor{JIN-1186-REGNAL-YEAR}{1186年·金以金世宗在位第25年作为诏令、赋役与外交文书的纪年框架}账本条目记录金在1186年的“金以金世宗在位第25年作为诏令、赋役与外交文书的纪年框架”。其背景是新岁颁历、年号和纪年是中枢与地方行政沟通的共同前提，后续影响为为同年政令、战事和地方材料提供日期锚点，不能单独推知具体政策执行。政权、军队或制度的变化须经由道路、文书、资源和地方中介才会落实，城市、乡村、边地和流动人群不可能在同一时间承受相同结果。

本条使用《金史》〈本纪〉、〈交聘表〉及相关志；《宋史》与《建炎以来系年要录》相关条，并标示“纪年框架可靠；该记录仅确证政权纪年连续，年内具体事项须由本年条另证”。这些材料能够钉住年序、命令、战役或局部地点，却不能直接推出人口、税收、身份和社会动机的总量。应区分诏令与实施、条约与边境生活、军报与伤亡统计，也不将官府的族群或行政称谓写成固定的社会本质。

从区域史角度看，该事件是资源、信息与权力重新连接的节点。不同地方会因市场、交通、宗教组织、家庭策略和环境条件而产生不同反应；材料的沉默限制我们对未被记录者所能作出的判断。

\eventanchor{JIN-1187-REGNAL-YEAR}{1187年·金以金世宗在位第26年作为诏令、赋役与外交文书的纪年框架}账本条目记录金在1187年的“金以金世宗在位第26年作为诏令、赋役与外交文书的纪年框架”。其背景是新岁颁历、年号和纪年是中枢与地方行政沟通的共同前提，后续影响为为同年政令、战事和地方材料提供日期锚点，不能单独推知具体政策执行。政权、军队或制度的变化须经由道路、文书、资源和地方中介才会落实，城市、乡村、边地和流动人群不可能在同一时间承受相同结果。

本条使用《金史》〈本纪〉、〈交聘表〉及相关志；《宋史》与《建炎以来系年要录》相关条，并标示“纪年框架可靠；该记录仅确证政权纪年连续，年内具体事项须由本年条另证”。这些材料能够钉住年序、命令、战役或局部地点，却不能直接推出人口、税收、身份和社会动机的总量。应区分诏令与实施、条约与边境生活、军报与伤亡统计，也不将官府的族群或行政称谓写成固定的社会本质。

从区域史角度看，该事件是资源、信息与权力重新连接的节点。不同地方会因市场、交通、宗教组织、家庭策略和环境条件而产生不同反应；材料的沉默限制我们对未被记录者所能作出的判断。

\eventanchor{JIN-1188-REGNAL-YEAR}{1188年·金以金世宗在位第27年作为诏令、赋役与外交文书的纪年框架}账本条目记录金在1188年的“金以金世宗在位第27年作为诏令、赋役与外交文书的纪年框架”。其背景是新岁颁历、年号和纪年是中枢与地方行政沟通的共同前提，后续影响为为同年政令、战事和地方材料提供日期锚点，不能单独推知具体政策执行。政权、军队或制度的变化须经由道路、文书、资源和地方中介才会落实，城市、乡村、边地和流动人群不可能在同一时间承受相同结果。

本条使用《金史》〈本纪〉、〈交聘表〉及相关志；《宋史》与《建炎以来系年要录》相关条，并标示“纪年框架可靠；该记录仅确证政权纪年连续，年内具体事项须由本年条另证”。这些材料能够钉住年序、命令、战役或局部地点，却不能直接推出人口、税收、身份和社会动机的总量。应区分诏令与实施、条约与边境生活、军报与伤亡统计，也不将官府的族群或行政称谓写成固定的社会本质。

从区域史角度看，该事件是资源、信息与权力重新连接的节点。不同地方会因市场、交通、宗教组织、家庭策略和环境条件而产生不同反应；材料的沉默限制我们对未被记录者所能作出的判断。

\eventanchor{JIN-1189-EVENT-099}{1189年·金世宗去世，章宗即位}账本条目记录金在1189年的“金世宗去世，章宗即位”。其背景是前期边防、财政和统治联盟的压力构成直接背景，具体条件须按本条材料分辨，后续影响为改变后续资源配置与政治选择；实际执行范围和社会影响仍须分项核验。政权、军队或制度的变化须经由道路、文书、资源和地方中介才会落实，城市、乡村、边地和流动人群不可能在同一时间承受相同结果。

本条使用《金史》〈本纪〉、〈交聘表〉及相关志；《宋史》与《建炎以来系年要录》相关条，并标示“编年主干可靠；细节、规模与动机须与对方材料、地方文书或考古材料互校”。这些材料能够钉住年序、命令、战役或局部地点，却不能直接推出人口、税收、身份和社会动机的总量。应区分诏令与实施、条约与边境生活、军报与伤亡统计，也不将官府的族群或行政称谓写成固定的社会本质。

从区域史角度看，该事件是资源、信息与权力重新连接的节点。不同地方会因市场、交通、宗教组织、家庭策略和环境条件而产生不同反应；材料的沉默限制我们对未被记录者所能作出的判断。

\eventanchor{JIN-1189-REGNAL-YEAR}{1189年·金以金世宗在位第28年作为诏令、赋役与外交文书的纪年框架}账本条目记录金在1189年的“金以金世宗在位第28年作为诏令、赋役与外交文书的纪年框架”。其背景是新岁颁历、年号和纪年是中枢与地方行政沟通的共同前提，后续影响为为同年政令、战事和地方材料提供日期锚点，不能单独推知具体政策执行。政权、军队或制度的变化须经由道路、文书、资源和地方中介才会落实，城市、乡村、边地和流动人群不可能在同一时间承受相同结果。

本条使用《金史》〈本纪〉、〈交聘表〉及相关志；《宋史》与《建炎以来系年要录》相关条，并标示“纪年框架可靠；该记录仅确证政权纪年连续，年内具体事项须由本年条另证”。这些材料能够钉住年序、命令、战役或局部地点，却不能直接推出人口、税收、身份和社会动机的总量。应区分诏令与实施、条约与边境生活、军报与伤亡统计，也不将官府的族群或行政称谓写成固定的社会本质。

从区域史角度看，该事件是资源、信息与权力重新连接的节点。不同地方会因市场、交通、宗教组织、家庭策略和环境条件而产生不同反应；材料的沉默限制我们对未被记录者所能作出的判断。

\eventanchor{JIN-1190-REGNAL-YEAR}{1190年·金以金章宗在位第1年作为诏令、赋役与外交文书的纪年框架}账本条目记录金在1190年的“金以金章宗在位第1年作为诏令、赋役与外交文书的纪年框架”。其背景是新岁颁历、年号和纪年是中枢与地方行政沟通的共同前提，后续影响为为同年政令、战事和地方材料提供日期锚点，不能单独推知具体政策执行。政权、军队或制度的变化须经由道路、文书、资源和地方中介才会落实，城市、乡村、边地和流动人群不可能在同一时间承受相同结果。

本条使用《金史》〈本纪〉、〈交聘表〉及相关志；《宋史》与《建炎以来系年要录》相关条，并标示“纪年框架可靠；该记录仅确证政权纪年连续，年内具体事项须由本年条另证”。这些材料能够钉住年序、命令、战役或局部地点，却不能直接推出人口、税收、身份和社会动机的总量。应区分诏令与实施、条约与边境生活、军报与伤亡统计，也不将官府的族群或行政称谓写成固定的社会本质。

从区域史角度看，该事件是资源、信息与权力重新连接的节点。不同地方会因市场、交通、宗教组织、家庭策略和环境条件而产生不同反应；材料的沉默限制我们对未被记录者所能作出的判断。

\eventanchor{JIN-1191-REGNAL-YEAR}{1191年·金以金章宗在位第2年作为诏令、赋役与外交文书的纪年框架}账本条目记录金在1191年的“金以金章宗在位第2年作为诏令、赋役与外交文书的纪年框架”。其背景是新岁颁历、年号和纪年是中枢与地方行政沟通的共同前提，后续影响为为同年政令、战事和地方材料提供日期锚点，不能单独推知具体政策执行。政权、军队或制度的变化须经由道路、文书、资源和地方中介才会落实，城市、乡村、边地和流动人群不可能在同一时间承受相同结果。

本条使用《金史》〈本纪〉、〈交聘表〉及相关志；《宋史》与《建炎以来系年要录》相关条，并标示“纪年框架可靠；该记录仅确证政权纪年连续，年内具体事项须由本年条另证”。这些材料能够钉住年序、命令、战役或局部地点，却不能直接推出人口、税收、身份和社会动机的总量。应区分诏令与实施、条约与边境生活、军报与伤亡统计，也不将官府的族群或行政称谓写成固定的社会本质。

从区域史角度看，该事件是资源、信息与权力重新连接的节点。不同地方会因市场、交通、宗教组织、家庭策略和环境条件而产生不同反应；材料的沉默限制我们对未被记录者所能作出的判断。

\eventanchor{JIN-1192-REGNAL-YEAR}{1192年·金以金章宗在位第3年作为诏令、赋役与外交文书的纪年框架}账本条目记录金在1192年的“金以金章宗在位第3年作为诏令、赋役与外交文书的纪年框架”。其背景是新岁颁历、年号和纪年是中枢与地方行政沟通的共同前提，后续影响为为同年政令、战事和地方材料提供日期锚点，不能单独推知具体政策执行。政权、军队或制度的变化须经由道路、文书、资源和地方中介才会落实，城市、乡村、边地和流动人群不可能在同一时间承受相同结果。

本条使用《金史》〈本纪〉、〈交聘表〉及相关志；《宋史》与《建炎以来系年要录》相关条，并标示“纪年框架可靠；该记录仅确证政权纪年连续，年内具体事项须由本年条另证”。这些材料能够钉住年序、命令、战役或局部地点，却不能直接推出人口、税收、身份和社会动机的总量。应区分诏令与实施、条约与边境生活、军报与伤亡统计，也不将官府的族群或行政称谓写成固定的社会本质。

从区域史角度看，该事件是资源、信息与权力重新连接的节点。不同地方会因市场、交通、宗教组织、家庭策略和环境条件而产生不同反应；材料的沉默限制我们对未被记录者所能作出的判断。

\eventanchor{JIN-1193-REGNAL-YEAR}{1193年·金以金章宗在位第4年作为诏令、赋役与外交文书的纪年框架}账本条目记录金在1193年的“金以金章宗在位第4年作为诏令、赋役与外交文书的纪年框架”。其背景是新岁颁历、年号和纪年是中枢与地方行政沟通的共同前提，后续影响为为同年政令、战事和地方材料提供日期锚点，不能单独推知具体政策执行。政权、军队或制度的变化须经由道路、文书、资源和地方中介才会落实，城市、乡村、边地和流动人群不可能在同一时间承受相同结果。

本条使用《金史》〈本纪〉、〈交聘表〉及相关志；《宋史》与《建炎以来系年要录》相关条，并标示“纪年框架可靠；该记录仅确证政权纪年连续，年内具体事项须由本年条另证”。这些材料能够钉住年序、命令、战役或局部地点，却不能直接推出人口、税收、身份和社会动机的总量。应区分诏令与实施、条约与边境生活、军报与伤亡统计，也不将官府的族群或行政称谓写成固定的社会本质。

从区域史角度看，该事件是资源、信息与权力重新连接的节点。不同地方会因市场、交通、宗教组织、家庭策略和环境条件而产生不同反应；材料的沉默限制我们对未被记录者所能作出的判断。

\eventanchor{JIN-1194-REGNAL-YEAR}{1194年·金以金章宗在位第5年作为诏令、赋役与外交文书的纪年框架}账本条目记录金在1194年的“金以金章宗在位第5年作为诏令、赋役与外交文书的纪年框架”。其背景是新岁颁历、年号和纪年是中枢与地方行政沟通的共同前提，后续影响为为同年政令、战事和地方材料提供日期锚点，不能单独推知具体政策执行。政权、军队或制度的变化须经由道路、文书、资源和地方中介才会落实，城市、乡村、边地和流动人群不可能在同一时间承受相同结果。

本条使用《金史》〈本纪〉、〈交聘表〉及相关志；《宋史》与《建炎以来系年要录》相关条，并标示“纪年框架可靠；该记录仅确证政权纪年连续，年内具体事项须由本年条另证”。这些材料能够钉住年序、命令、战役或局部地点，却不能直接推出人口、税收、身份和社会动机的总量。应区分诏令与实施、条约与边境生活、军报与伤亡统计，也不将官府的族群或行政称谓写成固定的社会本质。

从区域史角度看，该事件是资源、信息与权力重新连接的节点。不同地方会因市场、交通、宗教组织、家庭策略和环境条件而产生不同反应；材料的沉默限制我们对未被记录者所能作出的判断。

\eventanchor{JIN-1195-REGNAL-YEAR}{1195年·金以金章宗在位第6年作为诏令、赋役与外交文书的纪年框架}账本条目记录金在1195年的“金以金章宗在位第6年作为诏令、赋役与外交文书的纪年框架”。其背景是新岁颁历、年号和纪年是中枢与地方行政沟通的共同前提，后续影响为为同年政令、战事和地方材料提供日期锚点，不能单独推知具体政策执行。政权、军队或制度的变化须经由道路、文书、资源和地方中介才会落实，城市、乡村、边地和流动人群不可能在同一时间承受相同结果。

本条使用《金史》〈本纪〉、〈交聘表〉及相关志；《宋史》与《建炎以来系年要录》相关条，并标示“纪年框架可靠；该记录仅确证政权纪年连续，年内具体事项须由本年条另证”。这些材料能够钉住年序、命令、战役或局部地点，却不能直接推出人口、税收、身份和社会动机的总量。应区分诏令与实施、条约与边境生活、军报与伤亡统计，也不将官府的族群或行政称谓写成固定的社会本质。

从区域史角度看，该事件是资源、信息与权力重新连接的节点。不同地方会因市场、交通、宗教组织、家庭策略和环境条件而产生不同反应；材料的沉默限制我们对未被记录者所能作出的判断。

\eventanchor{JIN-1196-REGNAL-YEAR}{1196年·金以金章宗在位第7年作为诏令、赋役与外交文书的纪年框架}账本条目记录金在1196年的“金以金章宗在位第7年作为诏令、赋役与外交文书的纪年框架”。其背景是新岁颁历、年号和纪年是中枢与地方行政沟通的共同前提，后续影响为为同年政令、战事和地方材料提供日期锚点，不能单独推知具体政策执行。政权、军队或制度的变化须经由道路、文书、资源和地方中介才会落实，城市、乡村、边地和流动人群不可能在同一时间承受相同结果。

本条使用《金史》〈本纪〉、〈交聘表〉及相关志；《宋史》与《建炎以来系年要录》相关条，并标示“纪年框架可靠；该记录仅确证政权纪年连续，年内具体事项须由本年条另证”。这些材料能够钉住年序、命令、战役或局部地点，却不能直接推出人口、税收、身份和社会动机的总量。应区分诏令与实施、条约与边境生活、军报与伤亡统计，也不将官府的族群或行政称谓写成固定的社会本质。

从区域史角度看，该事件是资源、信息与权力重新连接的节点。不同地方会因市场、交通、宗教组织、家庭策略和环境条件而产生不同反应；材料的沉默限制我们对未被记录者所能作出的判断。

\eventanchor{JIN-1197-REGNAL-YEAR}{1197年·金以金章宗在位第8年作为诏令、赋役与外交文书的纪年框架}账本条目记录金在1197年的“金以金章宗在位第8年作为诏令、赋役与外交文书的纪年框架”。其背景是新岁颁历、年号和纪年是中枢与地方行政沟通的共同前提，后续影响为为同年政令、战事和地方材料提供日期锚点，不能单独推知具体政策执行。政权、军队或制度的变化须经由道路、文书、资源和地方中介才会落实，城市、乡村、边地和流动人群不可能在同一时间承受相同结果。

本条使用《金史》〈本纪〉、〈交聘表〉及相关志；《宋史》与《建炎以来系年要录》相关条，并标示“纪年框架可靠；该记录仅确证政权纪年连续，年内具体事项须由本年条另证”。这些材料能够钉住年序、命令、战役或局部地点，却不能直接推出人口、税收、身份和社会动机的总量。应区分诏令与实施、条约与边境生活、军报与伤亡统计，也不将官府的族群或行政称谓写成固定的社会本质。

从区域史角度看，该事件是资源、信息与权力重新连接的节点。不同地方会因市场、交通、宗教组织、家庭策略和环境条件而产生不同反应；材料的沉默限制我们对未被记录者所能作出的判断。

\eventanchor{JIN-1198-REGNAL-YEAR}{1198年·金以金章宗在位第9年作为诏令、赋役与外交文书的纪年框架}账本条目记录金在1198年的“金以金章宗在位第9年作为诏令、赋役与外交文书的纪年框架”。其背景是新岁颁历、年号和纪年是中枢与地方行政沟通的共同前提，后续影响为为同年政令、战事和地方材料提供日期锚点，不能单独推知具体政策执行。政权、军队或制度的变化须经由道路、文书、资源和地方中介才会落实，城市、乡村、边地和流动人群不可能在同一时间承受相同结果。

本条使用《金史》〈本纪〉、〈交聘表〉及相关志；《宋史》与《建炎以来系年要录》相关条，并标示“纪年框架可靠；该记录仅确证政权纪年连续，年内具体事项须由本年条另证”。这些材料能够钉住年序、命令、战役或局部地点，却不能直接推出人口、税收、身份和社会动机的总量。应区分诏令与实施、条约与边境生活、军报与伤亡统计，也不将官府的族群或行政称谓写成固定的社会本质。

从区域史角度看，该事件是资源、信息与权力重新连接的节点。不同地方会因市场、交通、宗教组织、家庭策略和环境条件而产生不同反应；材料的沉默限制我们对未被记录者所能作出的判断。

\eventanchor{JIN-1199-REGNAL-YEAR}{1199年·金以金章宗在位第10年作为诏令、赋役与外交文书的纪年框架}账本条目记录金在1199年的“金以金章宗在位第10年作为诏令、赋役与外交文书的纪年框架”。其背景是新岁颁历、年号和纪年是中枢与地方行政沟通的共同前提，后续影响为为同年政令、战事和地方材料提供日期锚点，不能单独推知具体政策执行。政权、军队或制度的变化须经由道路、文书、资源和地方中介才会落实，城市、乡村、边地和流动人群不可能在同一时间承受相同结果。

本条使用《金史》〈本纪〉、〈交聘表〉及相关志；《宋史》与《建炎以来系年要录》相关条，并标示“纪年框架可靠；该记录仅确证政权纪年连续，年内具体事项须由本年条另证”。这些材料能够钉住年序、命令、战役或局部地点，却不能直接推出人口、税收、身份和社会动机的总量。应区分诏令与实施、条约与边境生活、军报与伤亡统计，也不将官府的族群或行政称谓写成固定的社会本质。

从区域史角度看，该事件是资源、信息与权力重新连接的节点。不同地方会因市场、交通、宗教组织、家庭策略和环境条件而产生不同反应；材料的沉默限制我们对未被记录者所能作出的判断。

\eventanchor{JIN-1200-REGNAL-YEAR}{1200年·金以金章宗在位第11年作为诏令、赋役与外交文书的纪年框架}账本条目记录金在1200年的“金以金章宗在位第11年作为诏令、赋役与外交文书的纪年框架”。其背景是新岁颁历、年号和纪年是中枢与地方行政沟通的共同前提，后续影响为为同年政令、战事和地方材料提供日期锚点，不能单独推知具体政策执行。政权、军队或制度的变化须经由道路、文书、资源和地方中介才会落实，城市、乡村、边地和流动人群不可能在同一时间承受相同结果。

本条使用《金史》〈本纪〉、〈交聘表〉及相关志；《宋史》与《建炎以来系年要录》相关条，并标示“纪年框架可靠；该记录仅确证政权纪年连续，年内具体事项须由本年条另证”。这些材料能够钉住年序、命令、战役或局部地点，却不能直接推出人口、税收、身份和社会动机的总量。应区分诏令与实施、条约与边境生活、军报与伤亡统计，也不将官府的族群或行政称谓写成固定的社会本质。

从区域史角度看，该事件是资源、信息与权力重新连接的节点。不同地方会因市场、交通、宗教组织、家庭策略和环境条件而产生不同反应；材料的沉默限制我们对未被记录者所能作出的判断。

\eventanchor{JIN-1201-REGNAL-YEAR}{1201年·金以金章宗在位第12年作为诏令、赋役与外交文书的纪年框架}账本条目记录金在1201年的“金以金章宗在位第12年作为诏令、赋役与外交文书的纪年框架”。其背景是新岁颁历、年号和纪年是中枢与地方行政沟通的共同前提，后续影响为为同年政令、战事和地方材料提供日期锚点，不能单独推知具体政策执行。政权、军队或制度的变化须经由道路、文书、资源和地方中介才会落实，城市、乡村、边地和流动人群不可能在同一时间承受相同结果。

本条使用《金史》〈本纪〉、〈交聘表〉及相关志；《宋史》与《建炎以来系年要录》相关条，并标示“纪年框架可靠；该记录仅确证政权纪年连续，年内具体事项须由本年条另证”。这些材料能够钉住年序、命令、战役或局部地点，却不能直接推出人口、税收、身份和社会动机的总量。应区分诏令与实施、条约与边境生活、军报与伤亡统计，也不将官府的族群或行政称谓写成固定的社会本质。

从区域史角度看，该事件是资源、信息与权力重新连接的节点。不同地方会因市场、交通、宗教组织、家庭策略和环境条件而产生不同反应；材料的沉默限制我们对未被记录者所能作出的判断。

\eventanchor{JIN-1202-REGNAL-YEAR}{1202年·金以金章宗在位第13年作为诏令、赋役与外交文书的纪年框架}账本条目记录金在1202年的“金以金章宗在位第13年作为诏令、赋役与外交文书的纪年框架”。其背景是新岁颁历、年号和纪年是中枢与地方行政沟通的共同前提，后续影响为为同年政令、战事和地方材料提供日期锚点，不能单独推知具体政策执行。政权、军队或制度的变化须经由道路、文书、资源和地方中介才会落实，城市、乡村、边地和流动人群不可能在同一时间承受相同结果。

本条使用《金史》〈本纪〉、〈交聘表〉及相关志；《宋史》与《建炎以来系年要录》相关条，并标示“纪年框架可靠；该记录仅确证政权纪年连续，年内具体事项须由本年条另证”。这些材料能够钉住年序、命令、战役或局部地点，却不能直接推出人口、税收、身份和社会动机的总量。应区分诏令与实施、条约与边境生活、军报与伤亡统计，也不将官府的族群或行政称谓写成固定的社会本质。

从区域史角度看，该事件是资源、信息与权力重新连接的节点。不同地方会因市场、交通、宗教组织、家庭策略和环境条件而产生不同反应；材料的沉默限制我们对未被记录者所能作出的判断。

\subsection{制度、区域与材料边界}

\eventanchor{JIN-1203-REGNAL-YEAR}{1203年·金以金章宗在位第14年作为诏令、赋役与外交文书的纪年框架}账本条目记录金在1203年的“金以金章宗在位第14年作为诏令、赋役与外交文书的纪年框架”。其背景是新岁颁历、年号和纪年是中枢与地方行政沟通的共同前提，后续影响为为同年政令、战事和地方材料提供日期锚点，不能单独推知具体政策执行。政权、军队或制度的变化须经由道路、文书、资源和地方中介才会落实，城市、乡村、边地和流动人群不可能在同一时间承受相同结果。

本条使用《金史》〈本纪〉、〈交聘表〉及相关志；《宋史》与《建炎以来系年要录》相关条，并标示“纪年框架可靠；该记录仅确证政权纪年连续，年内具体事项须由本年条另证”。这些材料能够钉住年序、命令、战役或局部地点，却不能直接推出人口、税收、身份和社会动机的总量。应区分诏令与实施、条约与边境生活、军报与伤亡统计，也不将官府的族群或行政称谓写成固定的社会本质。

从区域史角度看，该事件是资源、信息与权力重新连接的节点。不同地方会因市场、交通、宗教组织、家庭策略和环境条件而产生不同反应；材料的沉默限制我们对未被记录者所能作出的判断。

\eventanchor{JIN-1204-REGNAL-YEAR}{1204年·金以金章宗在位第15年作为诏令、赋役与外交文书的纪年框架}账本条目记录金在1204年的“金以金章宗在位第15年作为诏令、赋役与外交文书的纪年框架”。其背景是新岁颁历、年号和纪年是中枢与地方行政沟通的共同前提，后续影响为为同年政令、战事和地方材料提供日期锚点，不能单独推知具体政策执行。政权、军队或制度的变化须经由道路、文书、资源和地方中介才会落实，城市、乡村、边地和流动人群不可能在同一时间承受相同结果。

本条使用《金史》〈本纪〉、〈交聘表〉及相关志；《宋史》与《建炎以来系年要录》相关条，并标示“纪年框架可靠；该记录仅确证政权纪年连续，年内具体事项须由本年条另证”。这些材料能够钉住年序、命令、战役或局部地点，却不能直接推出人口、税收、身份和社会动机的总量。应区分诏令与实施、条约与边境生活、军报与伤亡统计，也不将官府的族群或行政称谓写成固定的社会本质。

从区域史角度看，该事件是资源、信息与权力重新连接的节点。不同地方会因市场、交通、宗教组织、家庭策略和环境条件而产生不同反应；材料的沉默限制我们对未被记录者所能作出的判断。

\eventanchor{JIN-1205-REGNAL-YEAR}{1205年·金以金章宗在位第16年作为诏令、赋役与外交文书的纪年框架}账本条目记录金在1205年的“金以金章宗在位第16年作为诏令、赋役与外交文书的纪年框架”。其背景是新岁颁历、年号和纪年是中枢与地方行政沟通的共同前提，后续影响为为同年政令、战事和地方材料提供日期锚点，不能单独推知具体政策执行。政权、军队或制度的变化须经由道路、文书、资源和地方中介才会落实，城市、乡村、边地和流动人群不可能在同一时间承受相同结果。

本条使用《金史》〈本纪〉、〈交聘表〉及相关志；《宋史》与《建炎以来系年要录》相关条，并标示“纪年框架可靠；该记录仅确证政权纪年连续，年内具体事项须由本年条另证”。这些材料能够钉住年序、命令、战役或局部地点，却不能直接推出人口、税收、身份和社会动机的总量。应区分诏令与实施、条约与边境生活、军报与伤亡统计，也不将官府的族群或行政称谓写成固定的社会本质。

从区域史角度看，该事件是资源、信息与权力重新连接的节点。不同地方会因市场、交通、宗教组织、家庭策略和环境条件而产生不同反应；材料的沉默限制我们对未被记录者所能作出的判断。

\eventanchor{JIN-1206-EVENT-100}{1206年·金军进攻南宋，南宋北伐受挫}账本条目记录金与南宋在1206年的“金军进攻南宋，南宋北伐受挫”。其背景是前期边防、财政和统治联盟的压力构成直接背景，具体条件须按本条材料分辨，后续影响为改变后续资源配置与政治选择；实际执行范围和社会影响仍须分项核验。政权、军队或制度的变化须经由道路、文书、资源和地方中介才会落实，城市、乡村、边地和流动人群不可能在同一时间承受相同结果。

本条使用《金史》〈本纪〉、〈交聘表〉及相关志；《宋史》与《建炎以来系年要录》相关条，并标示“编年主干可靠；细节、规模与动机须与对方材料、地方文书或考古材料互校”。这些材料能够钉住年序、命令、战役或局部地点，却不能直接推出人口、税收、身份和社会动机的总量。应区分诏令与实施、条约与边境生活、军报与伤亡统计，也不将官府的族群或行政称谓写成固定的社会本质。

从区域史角度看，该事件是资源、信息与权力重新连接的节点。不同地方会因市场、交通、宗教组织、家庭策略和环境条件而产生不同反应；材料的沉默限制我们对未被记录者所能作出的判断。

\eventanchor{JIN-1206-REGNAL-YEAR}{1206年·金以金章宗在位第17年作为诏令、赋役与外交文书的纪年框架}账本条目记录金在1206年的“金以金章宗在位第17年作为诏令、赋役与外交文书的纪年框架”。其背景是新岁颁历、年号和纪年是中枢与地方行政沟通的共同前提，后续影响为为同年政令、战事和地方材料提供日期锚点，不能单独推知具体政策执行。政权、军队或制度的变化须经由道路、文书、资源和地方中介才会落实，城市、乡村、边地和流动人群不可能在同一时间承受相同结果。

本条使用《金史》〈本纪〉、〈交聘表〉及相关志；《宋史》与《建炎以来系年要录》相关条，并标示“纪年框架可靠；该记录仅确证政权纪年连续，年内具体事项须由本年条另证”。这些材料能够钉住年序、命令、战役或局部地点，却不能直接推出人口、税收、身份和社会动机的总量。应区分诏令与实施、条约与边境生活、军报与伤亡统计，也不将官府的族群或行政称谓写成固定的社会本质。

从区域史角度看，该事件是资源、信息与权力重新连接的节点。不同地方会因市场、交通、宗教组织、家庭策略和环境条件而产生不同反应；材料的沉默限制我们对未被记录者所能作出的判断。

\eventanchor{JIN-1207-REGNAL-YEAR}{1207年·金以金章宗在位第18年作为诏令、赋役与外交文书的纪年框架}账本条目记录金在1207年的“金以金章宗在位第18年作为诏令、赋役与外交文书的纪年框架”。其背景是新岁颁历、年号和纪年是中枢与地方行政沟通的共同前提，后续影响为为同年政令、战事和地方材料提供日期锚点，不能单独推知具体政策执行。政权、军队或制度的变化须经由道路、文书、资源和地方中介才会落实，城市、乡村、边地和流动人群不可能在同一时间承受相同结果。

本条使用《金史》〈本纪〉、〈交聘表〉及相关志；《宋史》与《建炎以来系年要录》相关条，并标示“纪年框架可靠；该记录仅确证政权纪年连续，年内具体事项须由本年条另证”。这些材料能够钉住年序、命令、战役或局部地点，却不能直接推出人口、税收、身份和社会动机的总量。应区分诏令与实施、条约与边境生活、军报与伤亡统计，也不将官府的族群或行政称谓写成固定的社会本质。

从区域史角度看，该事件是资源、信息与权力重新连接的节点。不同地方会因市场、交通、宗教组织、家庭策略和环境条件而产生不同反应；材料的沉默限制我们对未被记录者所能作出的判断。

\eventanchor{JIN-1208-EVENT-101}{1208年·金宋以嘉定和议结束战事}账本条目记录金与南宋在1208年的“金宋以嘉定和议结束战事”。其背景是前期边防、财政和统治联盟的压力构成直接背景，具体条件须按本条材料分辨，后续影响为改变后续资源配置与政治选择；实际执行范围和社会影响仍须分项核验。政权、军队或制度的变化须经由道路、文书、资源和地方中介才会落实，城市、乡村、边地和流动人群不可能在同一时间承受相同结果。

本条使用《金史》〈本纪〉、〈交聘表〉及相关志；《宋史》与《建炎以来系年要录》相关条，并标示“编年主干可靠；细节、规模与动机须与对方材料、地方文书或考古材料互校”。这些材料能够钉住年序、命令、战役或局部地点，却不能直接推出人口、税收、身份和社会动机的总量。应区分诏令与实施、条约与边境生活、军报与伤亡统计，也不将官府的族群或行政称谓写成固定的社会本质。

从区域史角度看，该事件是资源、信息与权力重新连接的节点。不同地方会因市场、交通、宗教组织、家庭策略和环境条件而产生不同反应；材料的沉默限制我们对未被记录者所能作出的判断。

\eventanchor{JIN-1208-REGNAL-YEAR}{1208年·金以金章宗在位第19年作为诏令、赋役与外交文书的纪年框架}账本条目记录金在1208年的“金以金章宗在位第19年作为诏令、赋役与外交文书的纪年框架”。其背景是新岁颁历、年号和纪年是中枢与地方行政沟通的共同前提，后续影响为为同年政令、战事和地方材料提供日期锚点，不能单独推知具体政策执行。政权、军队或制度的变化须经由道路、文书、资源和地方中介才会落实，城市、乡村、边地和流动人群不可能在同一时间承受相同结果。

本条使用《金史》〈本纪〉、〈交聘表〉及相关志；《宋史》与《建炎以来系年要录》相关条，并标示“纪年框架可靠；该记录仅确证政权纪年连续，年内具体事项须由本年条另证”。这些材料能够钉住年序、命令、战役或局部地点，却不能直接推出人口、税收、身份和社会动机的总量。应区分诏令与实施、条约与边境生活、军报与伤亡统计，也不将官府的族群或行政称谓写成固定的社会本质。

从区域史角度看，该事件是资源、信息与权力重新连接的节点。不同地方会因市场、交通、宗教组织、家庭策略和环境条件而产生不同反应；材料的沉默限制我们对未被记录者所能作出的判断。

\eventanchor{JIN-1209-REGNAL-YEAR}{1209年·金以卫绍王在位第1年作为诏令、赋役与外交文书的纪年框架}账本条目记录金在1209年的“金以卫绍王在位第1年作为诏令、赋役与外交文书的纪年框架”。其背景是新岁颁历、年号和纪年是中枢与地方行政沟通的共同前提，后续影响为为同年政令、战事和地方材料提供日期锚点，不能单独推知具体政策执行。政权、军队或制度的变化须经由道路、文书、资源和地方中介才会落实，城市、乡村、边地和流动人群不可能在同一时间承受相同结果。

本条使用《金史》〈本纪〉、〈交聘表〉及相关志；《宋史》与《建炎以来系年要录》相关条，并标示“纪年框架可靠；该记录仅确证政权纪年连续，年内具体事项须由本年条另证”。这些材料能够钉住年序、命令、战役或局部地点，却不能直接推出人口、税收、身份和社会动机的总量。应区分诏令与实施、条约与边境生活、军报与伤亡统计，也不将官府的族群或行政称谓写成固定的社会本质。

从区域史角度看，该事件是资源、信息与权力重新连接的节点。不同地方会因市场、交通、宗教组织、家庭策略和环境条件而产生不同反应；材料的沉默限制我们对未被记录者所能作出的判断。

\eventanchor{JIN-1210-REGNAL-YEAR}{1210年·金以卫绍王在位第2年作为诏令、赋役与外交文书的纪年框架}账本条目记录金在1210年的“金以卫绍王在位第2年作为诏令、赋役与外交文书的纪年框架”。其背景是新岁颁历、年号和纪年是中枢与地方行政沟通的共同前提，后续影响为为同年政令、战事和地方材料提供日期锚点，不能单独推知具体政策执行。政权、军队或制度的变化须经由道路、文书、资源和地方中介才会落实，城市、乡村、边地和流动人群不可能在同一时间承受相同结果。

本条使用《金史》〈本纪〉、〈交聘表〉及相关志；《宋史》与《建炎以来系年要录》相关条，并标示“纪年框架可靠；该记录仅确证政权纪年连续，年内具体事项须由本年条另证”。这些材料能够钉住年序、命令、战役或局部地点，却不能直接推出人口、税收、身份和社会动机的总量。应区分诏令与实施、条约与边境生活、军报与伤亡统计，也不将官府的族群或行政称谓写成固定的社会本质。

从区域史角度看，该事件是资源、信息与权力重新连接的节点。不同地方会因市场、交通、宗教组织、家庭策略和环境条件而产生不同反应；材料的沉默限制我们对未被记录者所能作出的判断。

\eventanchor{JIN-1211-EVENT-102}{1211年·蒙古大举攻金，金北方防线受冲击}账本条目记录金与蒙古在1211年的“蒙古大举攻金，金北方防线受冲击”。其背景是前期边防、财政和统治联盟的压力构成直接背景，具体条件须按本条材料分辨，后续影响为改变后续资源配置与政治选择；实际执行范围和社会影响仍须分项核验。政权、军队或制度的变化须经由道路、文书、资源和地方中介才会落实，城市、乡村、边地和流动人群不可能在同一时间承受相同结果。

本条使用《金史》〈本纪〉、〈交聘表〉及相关志；《宋史》与《建炎以来系年要录》相关条，并标示“编年主干可靠；细节、规模与动机须与对方材料、地方文书或考古材料互校”。这些材料能够钉住年序、命令、战役或局部地点，却不能直接推出人口、税收、身份和社会动机的总量。应区分诏令与实施、条约与边境生活、军报与伤亡统计，也不将官府的族群或行政称谓写成固定的社会本质。

从区域史角度看，该事件是资源、信息与权力重新连接的节点。不同地方会因市场、交通、宗教组织、家庭策略和环境条件而产生不同反应；材料的沉默限制我们对未被记录者所能作出的判断。

\eventanchor{JIN-1211-REGNAL-YEAR}{1211年·金以卫绍王在位第3年作为诏令、赋役与外交文书的纪年框架}账本条目记录金在1211年的“金以卫绍王在位第3年作为诏令、赋役与外交文书的纪年框架”。其背景是新岁颁历、年号和纪年是中枢与地方行政沟通的共同前提，后续影响为为同年政令、战事和地方材料提供日期锚点，不能单独推知具体政策执行。政权、军队或制度的变化须经由道路、文书、资源和地方中介才会落实，城市、乡村、边地和流动人群不可能在同一时间承受相同结果。

本条使用《金史》〈本纪〉、〈交聘表〉及相关志；《宋史》与《建炎以来系年要录》相关条，并标示“纪年框架可靠；该记录仅确证政权纪年连续，年内具体事项须由本年条另证”。这些材料能够钉住年序、命令、战役或局部地点，却不能直接推出人口、税收、身份和社会动机的总量。应区分诏令与实施、条约与边境生活、军报与伤亡统计，也不将官府的族群或行政称谓写成固定的社会本质。

从区域史角度看，该事件是资源、信息与权力重新连接的节点。不同地方会因市场、交通、宗教组织、家庭策略和环境条件而产生不同反应；材料的沉默限制我们对未被记录者所能作出的判断。

\eventanchor{JIN-1212-REGNAL-YEAR}{1212年·金以卫绍王在位第4年作为诏令、赋役与外交文书的纪年框架}账本条目记录金在1212年的“金以卫绍王在位第4年作为诏令、赋役与外交文书的纪年框架”。其背景是新岁颁历、年号和纪年是中枢与地方行政沟通的共同前提，后续影响为为同年政令、战事和地方材料提供日期锚点，不能单独推知具体政策执行。政权、军队或制度的变化须经由道路、文书、资源和地方中介才会落实，城市、乡村、边地和流动人群不可能在同一时间承受相同结果。

本条使用《金史》〈本纪〉、〈交聘表〉及相关志；《宋史》与《建炎以来系年要录》相关条，并标示“纪年框架可靠；该记录仅确证政权纪年连续，年内具体事项须由本年条另证”。这些材料能够钉住年序、命令、战役或局部地点，却不能直接推出人口、税收、身份和社会动机的总量。应区分诏令与实施、条约与边境生活、军报与伤亡统计，也不将官府的族群或行政称谓写成固定的社会本质。

从区域史角度看，该事件是资源、信息与权力重新连接的节点。不同地方会因市场、交通、宗教组织、家庭策略和环境条件而产生不同反应；材料的沉默限制我们对未被记录者所能作出的判断。

\eventanchor{JIN-1213-EVENT-103}{1213年·金宣宗即位并调整都城与军政策略}账本条目记录金在1213年的“金宣宗即位并调整都城与军政策略”。其背景是前期边防、财政和统治联盟的压力构成直接背景，具体条件须按本条材料分辨，后续影响为改变后续资源配置与政治选择；实际执行范围和社会影响仍须分项核验。政权、军队或制度的变化须经由道路、文书、资源和地方中介才会落实，城市、乡村、边地和流动人群不可能在同一时间承受相同结果。

本条使用《金史》〈本纪〉、〈交聘表〉及相关志；《宋史》与《建炎以来系年要录》相关条，并标示“编年主干可靠；细节、规模与动机须与对方材料、地方文书或考古材料互校”。这些材料能够钉住年序、命令、战役或局部地点，却不能直接推出人口、税收、身份和社会动机的总量。应区分诏令与实施、条约与边境生活、军报与伤亡统计，也不将官府的族群或行政称谓写成固定的社会本质。

从区域史角度看，该事件是资源、信息与权力重新连接的节点。不同地方会因市场、交通、宗教组织、家庭策略和环境条件而产生不同反应；材料的沉默限制我们对未被记录者所能作出的判断。

\eventanchor{JIN-1213-REGNAL-YEAR}{1213年·金以卫绍王在位第5年作为诏令、赋役与外交文书的纪年框架}账本条目记录金在1213年的“金以卫绍王在位第5年作为诏令、赋役与外交文书的纪年框架”。其背景是新岁颁历、年号和纪年是中枢与地方行政沟通的共同前提，后续影响为为同年政令、战事和地方材料提供日期锚点，不能单独推知具体政策执行。政权、军队或制度的变化须经由道路、文书、资源和地方中介才会落实，城市、乡村、边地和流动人群不可能在同一时间承受相同结果。

本条使用《金史》〈本纪〉、〈交聘表〉及相关志；《宋史》与《建炎以来系年要录》相关条，并标示“纪年框架可靠；该记录仅确证政权纪年连续，年内具体事项须由本年条另证”。这些材料能够钉住年序、命令、战役或局部地点，却不能直接推出人口、税收、身份和社会动机的总量。应区分诏令与实施、条约与边境生活、军报与伤亡统计，也不将官府的族群或行政称谓写成固定的社会本质。

从区域史角度看，该事件是资源、信息与权力重新连接的节点。不同地方会因市场、交通、宗教组织、家庭策略和环境条件而产生不同反应；材料的沉默限制我们对未被记录者所能作出的判断。

\eventanchor{JIN-1214-REGNAL-YEAR}{1214年·金以金宣宗在位第1年作为诏令、赋役与外交文书的纪年框架}账本条目记录金在1214年的“金以金宣宗在位第1年作为诏令、赋役与外交文书的纪年框架”。其背景是新岁颁历、年号和纪年是中枢与地方行政沟通的共同前提，后续影响为为同年政令、战事和地方材料提供日期锚点，不能单独推知具体政策执行。政权、军队或制度的变化须经由道路、文书、资源和地方中介才会落实，城市、乡村、边地和流动人群不可能在同一时间承受相同结果。

本条使用《金史》〈本纪〉、〈交聘表〉及相关志；《宋史》与《建炎以来系年要录》相关条，并标示“纪年框架可靠；该记录仅确证政权纪年连续，年内具体事项须由本年条另证”。这些材料能够钉住年序、命令、战役或局部地点，却不能直接推出人口、税收、身份和社会动机的总量。应区分诏令与实施、条约与边境生活、军报与伤亡统计，也不将官府的族群或行政称谓写成固定的社会本质。

从区域史角度看，该事件是资源、信息与权力重新连接的节点。不同地方会因市场、交通、宗教组织、家庭策略和环境条件而产生不同反应；材料的沉默限制我们对未被记录者所能作出的判断。

\eventanchor{JIN-1215-EVENT-104}{1215年·蒙古军攻取中都}账本条目记录金与蒙古在1215年的“蒙古军攻取中都”。其背景是前期边防、财政和统治联盟的压力构成直接背景，具体条件须按本条材料分辨，后续影响为改变后续资源配置与政治选择；实际执行范围和社会影响仍须分项核验。政权、军队或制度的变化须经由道路、文书、资源和地方中介才会落实，城市、乡村、边地和流动人群不可能在同一时间承受相同结果。

本条使用《金史》〈本纪〉、〈交聘表〉及相关志；《宋史》与《建炎以来系年要录》相关条，并标示“编年主干可靠；细节、规模与动机须与对方材料、地方文书或考古材料互校”。这些材料能够钉住年序、命令、战役或局部地点，却不能直接推出人口、税收、身份和社会动机的总量。应区分诏令与实施、条约与边境生活、军报与伤亡统计，也不将官府的族群或行政称谓写成固定的社会本质。

从区域史角度看，该事件是资源、信息与权力重新连接的节点。不同地方会因市场、交通、宗教组织、家庭策略和环境条件而产生不同反应；材料的沉默限制我们对未被记录者所能作出的判断。

\eventanchor{JIN-1215-REGNAL-YEAR}{1215年·金以金宣宗在位第2年作为诏令、赋役与外交文书的纪年框架}账本条目记录金在1215年的“金以金宣宗在位第2年作为诏令、赋役与外交文书的纪年框架”。其背景是新岁颁历、年号和纪年是中枢与地方行政沟通的共同前提，后续影响为为同年政令、战事和地方材料提供日期锚点，不能单独推知具体政策执行。政权、军队或制度的变化须经由道路、文书、资源和地方中介才会落实，城市、乡村、边地和流动人群不可能在同一时间承受相同结果。

本条使用《金史》〈本纪〉、〈交聘表〉及相关志；《宋史》与《建炎以来系年要录》相关条，并标示“纪年框架可靠；该记录仅确证政权纪年连续，年内具体事项须由本年条另证”。这些材料能够钉住年序、命令、战役或局部地点，却不能直接推出人口、税收、身份和社会动机的总量。应区分诏令与实施、条约与边境生活、军报与伤亡统计，也不将官府的族群或行政称谓写成固定的社会本质。

从区域史角度看，该事件是资源、信息与权力重新连接的节点。不同地方会因市场、交通、宗教组织、家庭策略和环境条件而产生不同反应；材料的沉默限制我们对未被记录者所能作出的判断。

\eventanchor{JIN-1216-REGNAL-YEAR}{1216年·金以金宣宗在位第3年作为诏令、赋役与外交文书的纪年框架}账本条目记录金在1216年的“金以金宣宗在位第3年作为诏令、赋役与外交文书的纪年框架”。其背景是新岁颁历、年号和纪年是中枢与地方行政沟通的共同前提，后续影响为为同年政令、战事和地方材料提供日期锚点，不能单独推知具体政策执行。政权、军队或制度的变化须经由道路、文书、资源和地方中介才会落实，城市、乡村、边地和流动人群不可能在同一时间承受相同结果。

本条使用《金史》〈本纪〉、〈交聘表〉及相关志；《宋史》与《建炎以来系年要录》相关条，并标示“纪年框架可靠；该记录仅确证政权纪年连续，年内具体事项须由本年条另证”。这些材料能够钉住年序、命令、战役或局部地点，却不能直接推出人口、税收、身份和社会动机的总量。应区分诏令与实施、条约与边境生活、军报与伤亡统计，也不将官府的族群或行政称谓写成固定的社会本质。

从区域史角度看，该事件是资源、信息与权力重新连接的节点。不同地方会因市场、交通、宗教组织、家庭策略和环境条件而产生不同反应；材料的沉默限制我们对未被记录者所能作出的判断。

\eventanchor{JIN-1217-EVENT-105}{1217年·金转而进攻南宋边境}账本条目记录金与南宋在1217年的“金转而进攻南宋边境”。其背景是前期边防、财政和统治联盟的压力构成直接背景，具体条件须按本条材料分辨，后续影响为改变后续资源配置与政治选择；实际执行范围和社会影响仍须分项核验。政权、军队或制度的变化须经由道路、文书、资源和地方中介才会落实，城市、乡村、边地和流动人群不可能在同一时间承受相同结果。

本条使用《金史》〈本纪〉、〈交聘表〉及相关志；《宋史》与《建炎以来系年要录》相关条，并标示“编年主干可靠；细节、规模与动机须与对方材料、地方文书或考古材料互校”。这些材料能够钉住年序、命令、战役或局部地点，却不能直接推出人口、税收、身份和社会动机的总量。应区分诏令与实施、条约与边境生活、军报与伤亡统计，也不将官府的族群或行政称谓写成固定的社会本质。

从区域史角度看，该事件是资源、信息与权力重新连接的节点。不同地方会因市场、交通、宗教组织、家庭策略和环境条件而产生不同反应；材料的沉默限制我们对未被记录者所能作出的判断。

\eventanchor{JIN-1217-REGNAL-YEAR}{1217年·金以金宣宗在位第4年作为诏令、赋役与外交文书的纪年框架}账本条目记录金在1217年的“金以金宣宗在位第4年作为诏令、赋役与外交文书的纪年框架”。其背景是新岁颁历、年号和纪年是中枢与地方行政沟通的共同前提，后续影响为为同年政令、战事和地方材料提供日期锚点，不能单独推知具体政策执行。政权、军队或制度的变化须经由道路、文书、资源和地方中介才会落实，城市、乡村、边地和流动人群不可能在同一时间承受相同结果。

本条使用《金史》〈本纪〉、〈交聘表〉及相关志；《宋史》与《建炎以来系年要录》相关条，并标示“纪年框架可靠；该记录仅确证政权纪年连续，年内具体事项须由本年条另证”。这些材料能够钉住年序、命令、战役或局部地点，却不能直接推出人口、税收、身份和社会动机的总量。应区分诏令与实施、条约与边境生活、军报与伤亡统计，也不将官府的族群或行政称谓写成固定的社会本质。

从区域史角度看，该事件是资源、信息与权力重新连接的节点。不同地方会因市场、交通、宗教组织、家庭策略和环境条件而产生不同反应；材料的沉默限制我们对未被记录者所能作出的判断。

\eventanchor{JIN-1218-REGNAL-YEAR}{1218年·金以金宣宗在位第5年作为诏令、赋役与外交文书的纪年框架}账本条目记录金在1218年的“金以金宣宗在位第5年作为诏令、赋役与外交文书的纪年框架”。其背景是新岁颁历、年号和纪年是中枢与地方行政沟通的共同前提，后续影响为为同年政令、战事和地方材料提供日期锚点，不能单独推知具体政策执行。政权、军队或制度的变化须经由道路、文书、资源和地方中介才会落实，城市、乡村、边地和流动人群不可能在同一时间承受相同结果。

本条使用《金史》〈本纪〉、〈交聘表〉及相关志；《宋史》与《建炎以来系年要录》相关条，并标示“纪年框架可靠；该记录仅确证政权纪年连续，年内具体事项须由本年条另证”。这些材料能够钉住年序、命令、战役或局部地点，却不能直接推出人口、税收、身份和社会动机的总量。应区分诏令与实施、条约与边境生活、军报与伤亡统计，也不将官府的族群或行政称谓写成固定的社会本质。

从区域史角度看，该事件是资源、信息与权力重新连接的节点。不同地方会因市场、交通、宗教组织、家庭策略和环境条件而产生不同反应；材料的沉默限制我们对未被记录者所能作出的判断。

\eventanchor{JIN-1219-REGNAL-YEAR}{1219年·金以金宣宗在位第6年作为诏令、赋役与外交文书的纪年框架}账本条目记录金在1219年的“金以金宣宗在位第6年作为诏令、赋役与外交文书的纪年框架”。其背景是新岁颁历、年号和纪年是中枢与地方行政沟通的共同前提，后续影响为为同年政令、战事和地方材料提供日期锚点，不能单独推知具体政策执行。政权、军队或制度的变化须经由道路、文书、资源和地方中介才会落实，城市、乡村、边地和流动人群不可能在同一时间承受相同结果。

本条使用《金史》〈本纪〉、〈交聘表〉及相关志；《宋史》与《建炎以来系年要录》相关条，并标示“纪年框架可靠；该记录仅确证政权纪年连续，年内具体事项须由本年条另证”。这些材料能够钉住年序、命令、战役或局部地点，却不能直接推出人口、税收、身份和社会动机的总量。应区分诏令与实施、条约与边境生活、军报与伤亡统计，也不将官府的族群或行政称谓写成固定的社会本质。

从区域史角度看，该事件是资源、信息与权力重新连接的节点。不同地方会因市场、交通、宗教组织、家庭策略和环境条件而产生不同反应；材料的沉默限制我们对未被记录者所能作出的判断。

\eventanchor{JIN-1220-REGNAL-YEAR}{1220年·金以金宣宗在位第7年作为诏令、赋役与外交文书的纪年框架}账本条目记录金在1220年的“金以金宣宗在位第7年作为诏令、赋役与外交文书的纪年框架”。其背景是新岁颁历、年号和纪年是中枢与地方行政沟通的共同前提，后续影响为为同年政令、战事和地方材料提供日期锚点，不能单独推知具体政策执行。政权、军队或制度的变化须经由道路、文书、资源和地方中介才会落实，城市、乡村、边地和流动人群不可能在同一时间承受相同结果。

本条使用《金史》〈本纪〉、〈交聘表〉及相关志；《宋史》与《建炎以来系年要录》相关条，并标示“纪年框架可靠；该记录仅确证政权纪年连续，年内具体事项须由本年条另证”。这些材料能够钉住年序、命令、战役或局部地点，却不能直接推出人口、税收、身份和社会动机的总量。应区分诏令与实施、条约与边境生活、军报与伤亡统计，也不将官府的族群或行政称谓写成固定的社会本质。

从区域史角度看，该事件是资源、信息与权力重新连接的节点。不同地方会因市场、交通、宗教组织、家庭策略和环境条件而产生不同反应；材料的沉默限制我们对未被记录者所能作出的判断。

\eventanchor{JIN-1221-REGNAL-YEAR}{1221年·金以金宣宗在位第8年作为诏令、赋役与外交文书的纪年框架}账本条目记录金在1221年的“金以金宣宗在位第8年作为诏令、赋役与外交文书的纪年框架”。其背景是新岁颁历、年号和纪年是中枢与地方行政沟通的共同前提，后续影响为为同年政令、战事和地方材料提供日期锚点，不能单独推知具体政策执行。政权、军队或制度的变化须经由道路、文书、资源和地方中介才会落实，城市、乡村、边地和流动人群不可能在同一时间承受相同结果。

本条使用《金史》〈本纪〉、〈交聘表〉及相关志；《宋史》与《建炎以来系年要录》相关条，并标示“纪年框架可靠；该记录仅确证政权纪年连续，年内具体事项须由本年条另证”。这些材料能够钉住年序、命令、战役或局部地点，却不能直接推出人口、税收、身份和社会动机的总量。应区分诏令与实施、条约与边境生活、军报与伤亡统计，也不将官府的族群或行政称谓写成固定的社会本质。

从区域史角度看，该事件是资源、信息与权力重新连接的节点。不同地方会因市场、交通、宗教组织、家庭策略和环境条件而产生不同反应；材料的沉默限制我们对未被记录者所能作出的判断。
